\clearpage
\SourceLeaf{106}\pdfbookmark[1]{Abschnitt III.}{sectio-tertia}
\begin{center}
{\large ABSCHNITT III.}\par\medskip
\emph{Beziehungen zwischen mehreren Orten auf der Bahn.}\par\bigskip
\hypertarget{native-106-article078}{}\pdfbookmark[2]{78.}{article-078}78.
\end{center}
Der Vergleich zweier oder mehrerer Orte eines Himmelskörpers sowohl auf seiner Bahn als auch im Raum liefert so viele elegante Sätze, dass sie leicht einen ganzen Band füllen könnten. Eine erschöpfende Behandlung dieses reichen Gegenstands liegt außerhalb unseres Vorhabens. Unser Ziel ist es vor allem, daraus umfangreiche Hilfsmittel zur Lösung der wichtigen Aufgabe zu gewinnen, unbekannte Bahnen aus Beobachtungen zu bestimmen. Wir lassen daher beiseite, was zu weit von unserem Vorhaben abliegt, und entwickeln umso sorgfältiger alles, was in irgendeiner Weise dazu beitragen kann. Der eigentlichen Untersuchung stellen wir einige trigonometrische Sätze voran, die häufig gebraucht werden und auf die wir wiederholt zurückgreifen müssen.
\par
I. Es seien $A$, $B$ und $C$ beliebige Winkel. Dann gilt
\[
\begin{aligned}
\sin A\sin(C-B)+\sin B\sin(A-C)+\sin C\sin(B-A)&=0\\
\cos A\sin(C-B)+\cos B\sin(A-C)+\cos C\sin(B-A)&=0.
\end{aligned}
\]
\par
II. Sollen zwei Größen $p$ und $P$ aus Gleichungen der Form
\[
\begin{aligned}
p\sin(A-P)&=a\\
p\sin(B-P)&=b
\end{aligned}
\]
bestimmt werden, so geschieht dies allgemein mithilfe der Formeln
\SourceLeaf{107}
\[
\begin{aligned}
p\sin(B-A)\sin(H-P)&=b\sin(H-A)-a\sin(H-B)\\
p\sin(B-A)\cos(H-P)&=b\cos(H-A)-a\cos(H-B),
\end{aligned}
\]
wobei $H$ ein beliebiger Winkel ist. Daraus erhält man (Artikel 14, II) den Winkel $H-P$ sowie $p\sin(B-A)$ und aus diesen wiederum $P$ und $p$. Meist kommt die Bedingung hinzu, dass $p$ positiv sein muss. Dadurch wird die Mehrdeutigkeit bei der Bestimmung des Winkels $H-P$ aus seinem Tangens beseitigt. Fehlt diese Bedingung, so steht die Wahl frei. Um die Rechnung möglichst bequem zu machen, setzt man den beliebigen Winkel $H$ entweder $=A$, $=B$ oder $=\tfrac{1}{2}(A+B)$. Im ersten Fall lauten die Gleichungen zur Bestimmung von $P$ und $p$
\[
\begin{aligned}
p\sin(A-P)&=a\\
p\cos(A-P)&=\frac{b-a\cos(B-A)}{\sin(B-A)}\mathpunct{\cdot}
\end{aligned}
\]
Die Gleichungen im zweiten Fall sind ganz analog. Im dritten Fall lauten sie
\[
\begin{aligned}
p\sin(\tfrac{1}{2}A+\tfrac{1}{2}B-P)&=\frac{b+a}{2\cos\tfrac{1}{2}(B-A)}\\
p\cos(\tfrac{1}{2}A+\tfrac{1}{2}B-P)&=\frac{b-a}{2\sin\tfrac{1}{2}(B-A)}\mathpunct{\cdot}
\end{aligned}
\]
\par
Führt man also den Hilfswinkel $\zeta$ mit dem Tangens $=\frac{a}{b}$ ein, so erhält man $P$ nach der Formel
\[
\operatorname{tang}(\tfrac{1}{2}A+\tfrac{1}{2}B-P)=\operatorname{tang}(45^0+\zeta)\operatorname{tang}\tfrac{1}{2}(B-A)
\]
und anschließend $p$ nach einer der vorangehenden Formeln, wobei gilt
\[
\begin{aligned}
\tfrac{1}{2}(b+a)&=\sin(45^0+\zeta)\BareRoot{\frac{ab}{\sin2\zeta}}\\
&=\frac{a\sin(45^0+\zeta)}{\sin\zeta\BareRoot{2}}=\frac{b\sin(45^0+\zeta)}{\cos\zeta\BareRoot{2}}\\[.4em]
\tfrac{1}{2}(b-a)&=\cos(45^0+\zeta)\BareRoot{\frac{ab}{\sin2\zeta}}\\
&=\frac{a\cos(45^0+\zeta)}{\sin\zeta\BareRoot{2}}=\frac{b\cos(45^0+\zeta)}{\cos\zeta\BareRoot{2}}\mathpunct{\cdot}
\end{aligned}
\]
\par
III. Sollen $p$ und $P$ aus den Gleichungen
\[
\begin{aligned}
p\cos(A-P)&=a\\
p\cos(B-P)&=b,
\end{aligned}
\]
bestimmt werden, so ließe sich alles unter II Erläuterte unmittelbar anwenden, wenn man dort überall für $A$ und $B$ die Ausdrücke $90^0+A$ und $90^0+B$ schriebe. Der bequemeren Anwendung halber geben wir die daraus folgenden Formeln ausdrücklich an. Die allgemeinen Formeln lauten
\SourceLeaf{108}
\[
\begin{aligned}
p\sin(B-A)\sin(H-P)&=-b\cos(H-A)+a\cos(H-B)\\
p\sin(B-A)\cos(H-P)&=b\sin(H-A)-a\sin(H-B).
\end{aligned}
\]
Für $H=A$ gehen sie somit über in
\[
\begin{aligned}
p\sin(A-P)&=\frac{a\cos(B-A)-b}{\sin(B-A)}\\
p\cos(A-P)&=a.
\end{aligned}
\]
Für $H=B$ nehmen sie eine ähnliche Form an. Für $H=\tfrac{1}{2}(A+B)$ lauten sie
\[
\begin{aligned}
p\sin(\tfrac{1}{2}A+\tfrac{1}{2}B-P)&=\frac{a-b}{2\sin\tfrac{1}{2}(B-A)}\\
p\cos(\tfrac{1}{2}A+\tfrac{1}{2}B-P)&=\frac{a+b}{2\cos\tfrac{1}{2}(B-A)}\mathpunct{\raisebox{.45ex}{,}}
\end{aligned}
\]
sodass sich mit dem Hilfswinkel $\zeta$, dessen Tangens $=\frac{a}{b}$ ist, ergibt
\[
\operatorname{cotang}(\tfrac{1}{2}A+\tfrac{1}{2}B-P)=\operatorname{cotang}(\zeta-45^0)\operatorname{tang}\tfrac{1}{2}(B-A).
\]
\par
Will man ferner $p$ unmittelbar aus $a$ und $b$ bestimmen, ohne zuvor den Winkel $P$ zu berechnen, so gilt die Formel
\[
p\sin(B-A)=\BareRoot{(aa+bb-2ab\cos(B-A))}
\]
sowohl für die vorliegende Aufgabe als auch für die unter II.
